\chapter{第一单元}

\Needspace{12\baselineskip}

\section{第一单元素养评价（A）}

\Needspace{8\baselineskip}

\pdfbookmark[2]{原书第1页}{orig-1}

\subsection*{原书第1页}

% p001_s01
\BookQuestion{一、1}
添加括号使\(671+29\times4\div2\)的运算顺序变为先求和，再求积，最后求商，则原式变为（\hspace{1.8em}）。\par
\begin{BookAnswer}
先加括号把\(671+29\)作为整体，再乘\(4\)，最后除以\(2\)。应填\((671+29)\times4\div2\)。
\end{BookAnswer}

% p001_s02
\BookQuestion{一、2（1）}
先在方框里填上合适的数，再列出综合算式。（1）见运算流程图。综合算式：\BookBlank{}。\par
\BookFigure{figures/p001_s02-01.tex}
\begin{BookAnswer}
顺箭头依次求数：\(136-46=90\)，\allowbreak{}\(450\div90=5\)，\allowbreak{}\(5\times22=110\)。三格依次为\(90,5,110\)；综合算式为\(450\div(136-46)\times22\)。
\BookFigure{figures/p001_s02-02.tex}
\end{BookAnswer}

% p001_s03
\BookQuestion{一、2（2）}
先在方框里填上合适的数，再列出综合算式。（2）见运算流程图。综合算式：\BookBlank{}。\par
\BookFigure{figures/p001_s03-01.tex}
\begin{BookAnswer}
顺箭头依次求数：\(252\div42=6\)，\allowbreak{}\(219-6=213\)，\allowbreak{}\(731-213=518\)。三格依次为\(6,213,518\)；综合算式为\(731-(219-252\div42)\)。
\BookFigure{figures/p001_s03-02.tex}
\end{BookAnswer}

% p001_s04
\BookQuestion{一、3}
根据运算顺序，把\(53-35=18\)，\allowbreak{}\(18\times4=72\)，\allowbreak{}\(720\div72=10\)改写成综合算式是（\hspace{1.8em}）。\par
\begin{BookAnswer}
先算\(53-35\)，所得差乘\(4\)，最后用\(720\)除以这个积。应填\(720\div[(53-35)\times4]\)，得数为\(10\)。
\end{BookAnswer}

% p001_s05
\BookQuestion{一、4}
要使\(8\times(\square\div5)=120\)，方框里应该填（\hspace{1.8em}）。\par
\begin{BookAnswer}
逆着运算顺序还原：括号内的商是\(120\div8=15\)，被除数是\(15\times5=75\)。填\(75\)；检验\(8\times(75\div5)=120\)。
\end{BookAnswer}

% p001_s06
\BookQuestion{一、5}
小丽在计算一道除法算式时，把除数32看成了23，算出来的商是34，余数是18。这道计算题正确的结果是（\hspace{1.8em}）。\par
\begin{BookAnswer}
看错除数时，被除数并没有改变。先用“除数\(\times\)商\(+\)余数”还原被除数，再除以\(32\)。\[\begin{aligned}&23\times34+18\\&=782+18\\&=800\\&800\div32=25\end{aligned}\]所以正确结果是\(25\)。
\end{BookAnswer}

% p001_s07
\BookQuestion{一、6（1）}
在圆圈里填上“\(>\)”“\(<\)”或“\(=\)”：\(122\times40-15+328\ \bigcirc\ 122\times(40-15)+328\)。\par
\begin{BookAnswer}
左式按先乘后加减算：\(122\times40-15+328=5193\)；右式先算括号：\(122\times(40-15)+328=3378\)。\(5193>3378\)，填\(>\)。
\end{BookAnswer}

% p001_s08
\BookQuestion{一、6（2）}
在圆圈里填上“\(>\)”“\(<\)”或“\(=\)”：\(560\div(14-6)\times5\ \bigcirc\ 560\div[(14-6)\times5]\)。\par
\begin{BookAnswer}
\(14-6=8\)。左式为\(560\div8\times5=350\)，右式为\(560\div(8\times5)=14\)。\(350>14\)，填\(>\)。
\end{BookAnswer}

% p001_s09
\BookQuestion{一、7（跨页）}
傍晚，王叔叔在小区里走路锻炼，步行2分钟后，他手机计步APP显示如图。照这样计算，后来他又步行32分钟。\(130+130\div2\times32\)算式求的问题是（\hspace{1.8em}）。（题干跨至原书第2页。）\par
\par\smallskip\Needspace{4\baselineskip}\begingroup\small\setlength{\tabcolsep}{3pt}
\noindent\begin{tabularx}{\linewidth}{|*{3}{>{\centering\arraybackslash}X|}}\hline
步数 & 步行距离 & 消耗热量\\\hline
200步 & 130米 & 4千卡\\\hline
\end{tabularx}\par
\endgroup\smallskip
\begin{BookAnswer}
\(130\div2\)表示每分钟走的米数，再乘\(32\)求后来走的路程，最后加上最初的\(130\)米。因此算式求“王叔叔一共走了多少米”。完整计算见下一页。
\end{BookAnswer}

\Needspace{8\baselineskip}

\pdfbookmark[2]{原书第2页}{orig-2}

\subsection*{原书第2页}

% p002_s01
\BookQuestion{一、7（续）}
（接原书第1页第7题）王叔叔步行2分钟走130米，后来又步行32分钟。\(130+130\div2\times32\)算式求的问题是（\hspace{1.8em}）。\par
\begin{BookAnswer}
“照这样计算”表示速度保持不变。先由前一段求\(130\div2=65\text{(米/分)}\)，再用“路程\(=\)速度\(\times\)时间”分别求后段或全程。\par 应填“一共走了多少米”。\[\begin{aligned}&130+130\div2\times32\\&=130+65\times32\\&=130+2080\\&=2210\text{(米)}\end{aligned}\]\par \par\noindent\begin{minipage}{\linewidth}\textbf{另一方法：先合并时间}\par \[\begin{aligned}&130\div2\times(2+32)\\&=65\times34\\&=2210\text{(米)}\end{aligned}\]\end{minipage}\par \par\textbf{简便提示（学完运算律后）}\par 先求速度的主法保留。计算后段路程时，可用整数乘除的关系先算32除以2；若先合并总时间34分钟，也可先算34除以2。只调整数值计算的次序，没有把32分钟改成总时间。\[\begin{aligned}&130+130\div2\times32\\={}&130+130\times(32\div2)\\={}&130+130\times16\\={}&130+2080\\={}&2210\text{(米)}\end{aligned}\]先合并时间的一路可简算为：\[\begin{aligned}&130\div2\times(2+32)\\={}&130\div2\times34\\={}&130\times(34\div2)\\={}&130\times17\\={}&2210\text{(米)}\end{aligned}\]
\end{BookAnswer}

% p002_s02
\BookQuestion{一、8}
学校手工社团买来320米长的绳子。做手工时，先剪了15段，每段长6米，又剪了48米备用，还剩下（\hspace{1.8em}）米绳子。梳理条件信息，应先求（\hspace{1.8em}），再求（\hspace{1.8em}），最后求（\hspace{1.8em}）。列综合算式为：（\hspace{1.8em}）。\par
\begin{BookAnswer}
三个文字空依次可填：\textbf{做手工用去的米数、共用去的米数、剩下的米数}；数值空填\(182\)。综合算式可填\(320-(6\times15+48)\)。\par \par\noindent\begin{minipage}{\linewidth}\textbf{方法一：先合计用去的长度}\par \[\begin{aligned}&320-(6\times15+48)\\&=320-(90+48)\\&=320-138\\&=182\text{(米)}\end{aligned}\]答：还剩\(182\)米绳子。\end{minipage}\par \par \par\noindent\begin{minipage}{\linewidth}\textbf{方法二：依剪绳的次序扣除}\par \[\begin{aligned}&320-6\times15-48\\&=320-90-48\\&=230-48\\&=182\text{(米)}\end{aligned}\]答：还剩\(182\)米绳子。\end{minipage}\par \par 回顾检验：\(182+6\times15+48=320\)，即“剩下的\(+\)手工用去的\(+\)备用的\(=\)总长”，与题意相符。
\end{BookAnswer}

% p002_s03
\BookQuestion{二、1}
计算\(4\times25\div25\times4\)的结果，正确的是（\hspace{1.8em}）。\par
\noindent A.1\par
\noindent B.16\par
\noindent C.0\par
\noindent D.都不对\par
\begin{BookAnswer}
选\(\mathrm B\)。乘除同级，从左往右：\(4\times25\div25\times4=100\div25\times4=16\)。也可把\(\times25\div25\)抵消，剩\(4\times4=16\)。
\end{BookAnswer}

% p002_s04
\BookQuestion{二、2}
计算\(21\times[950\div(22+28)-12]\)时，运算顺序是（\hspace{1.8em}）。\par
\noindent A.加→减→乘→除\par
\noindent B.除→加→减→乘\par
\noindent C.加→除→减→乘\par
\noindent D.加→除→乘→减\par
\begin{BookAnswer}
选\(\mathrm C\)。先算\(22+28\)，再算\(950\div50\)，接着减\(12\)，最后乘\(21\)；顺序为“加、除、减、乘”。
\end{BookAnswer}

% p002_s05
\BookQuestion{二、3}
如果\(90\times35-25\div5\)先算减法，再算除法，最后算乘法，正确的算式是（\hspace{1.8em}）。\par
\noindent A.\(90\times(35-25)\div5\)\par
\noindent B.\(90\times[(35-25)\div5]\)\par
\noindent C.\(90\times(35-25\div5)\)\par
\noindent D.\((90\times35-25)\div5\)\par
\begin{BookAnswer}
选\(\mathrm B\)。先减得到\(35-25\)，再将这个差除以\(5\)，最后乘\(90\)，应写\(90\times[(35-25)\div5]\)。
\end{BookAnswer}

% p002_s06
\BookQuestion{二、4}
豆豆不小心将作业本弄脏了，遮住了解决问题的部分内容。已知爸爸往物业费账户里充值了945元，第一季度扣除了405元。豆豆列式\((945-405)\div(405\div3)=4\)（个）。本题想要解决的问题是（\hspace{1.8em}）。\par
\noindent A.物业费账户里还剩多少元\par
\noindent B.一共需要缴纳多少元物业费\par
\noindent C.物业费账户里的钱够扣除几个月\par
\noindent D.物业费账户里剩余的钱够扣除几个月\par
\begin{BookAnswer}
选\(\mathrm D\)。\(945-405\)是账户余额，\allowbreak{}\(405\div3\)是每月物业费；“余额\(\div\)每月费用”求剩余的钱够交几个月。
\end{BookAnswer}

% p002_s07
\BookQuestion{二、5}
爸爸带豆豆去离家2000米的图书馆看书，父子俩步行速度是70米/分，他们从家出发25分钟后，位置大约在（\hspace{1.8em}）处。\par
\noindent A.A处\par
\noindent B.B处\par
\noindent C.C处\par
\noindent D.D处\par
\BookFigure{figures/p002_s07-01.tex}
\begin{BookAnswer}
选\(\mathrm C\)。父子\(25\)分钟走\(70\times25=1750\)米，距图书馆还剩\(2000-1750=250\)米，即全程的八分之一。图中与此位置最接近的是\(C\)。
\end{BookAnswer}

% p002_s08
\BookQuestion{三、1}
直接写得数：\(84\div4=\ \)，\allowbreak{}\(23\times3=\ \)，\allowbreak{}\(360\times4=\ \)，\allowbreak{}\(8+36\div4+8=\ \)，\allowbreak{}\(163-45=\ \)，\allowbreak{}\(98\div7=\ \)，\allowbreak{}\(84\div6=\ \)，\allowbreak{}\((28+32)\div5=\ \)。\par
\begin{BookAnswer}
按题目从左到右、先上后下，得数依次为\(21,69,1440,25\)；\allowbreak{}\(118,14,14,12\)。\par 其中\(8+36\div4+8=8+9+8=25\)，\allowbreak{}\((28+32)\div5=60\div5=12\)；其余直接做整数乘除或减法。
\end{BookAnswer}

\Needspace{8\baselineskip}

\pdfbookmark[2]{原书第3页}{orig-3}

\subsection*{原书第3页}

% p003_s01
\BookQuestion{三、2（1）}
计算下列各题：\((432\div36+45)\times24\)。\par
\begin{BookAnswer}
\[\begin{aligned}&(432\div36+45)\times24\\&=(12+45)\times24\\&=57\times24\\&=1368\end{aligned}\]
\end{BookAnswer}

% p003_s02
\BookQuestion{三、2（2）}
计算下列各题：\(91-(352+272)\div39\)。\par
\begin{BookAnswer}
\[\begin{aligned}&91-(352+272)\div39\\&=91-624\div39\\&=91-16\\&=75\end{aligned}\]
\end{BookAnswer}

% p003_s03
\BookQuestion{三、2（3）}
计算下列各题：\(592-(118\times4-215)\)。\par
\begin{BookAnswer}
\[\begin{aligned}&592-(118\times4-215)\\&=592-(472-215)\\&=592-257\\&=335\end{aligned}\]\par\textbf{简便提示（学完运算律后）}\par \[\begin{aligned}&592-(472-215)\\={}&592-472+215\\={}&120+215\\={}&335\end{aligned}\]
\end{BookAnswer}

% p003_s04
\BookQuestion{三、2（4）}
计算下列各题：\(76\times[215-(62+78)]\)。\par
\begin{BookAnswer}
\[\begin{aligned}&76\times[215-(62+78)]\\&=76\times[215-140]\\&=76\times75\\&=5700\end{aligned}\]\par\textbf{简便提示（学完运算律后）}\par \[\begin{aligned}&76\times75\\={}&19\times4\times75\\={}&19\times(4\times75)\\={}&19\times300\\={}&5700\end{aligned}\]
\end{BookAnswer}

% p003_s05
\BookQuestion{三、2（5）}
计算下列各题：\([396\div(302-291)]\times58\)。\par
\begin{BookAnswer}
\[\begin{aligned}&[396\div(302-291)]\times58\\&=[396\div11]\times58\\&=36\times58\\&=2088\end{aligned}\]\par\textbf{简便提示（学完运算律后）}\par 前面的括号和除法仍按原顺序求得36。最后把58看成60少2，先多算两份36再减回；两个乘法在同一步完成。\[\begin{aligned}&36\times58\\={}&36\times(60-2)\\={}&36\times60-36\times2\\={}&2160-72\\={}&2088\end{aligned}\]
\end{BookAnswer}

% p003_s06
\BookQuestion{三、2（6）}
计算下列各题：\(184\div[560\div(328-258)]\)。\par
\begin{BookAnswer}
\[\begin{aligned}&184\div[560\div(328-258)]\\&=184\div[560\div70]\\&=184\div8\\&=23\end{aligned}\]
\end{BookAnswer}

% p003_s07
\BookQuestion{四、1}
四年级18个班在“环保回收·废纸利用”活动中，平均每班回收废纸4千克。学校回收站已回收处理了32千克。剩下的废纸要统一分4批送去再生纸工厂。平均每批送多少千克？先分别从条件、从问题开始分析，填好流程图，再列式解答。\par
\BookFigure{figures/p003_s07-01.tex}
\BookFigure{figures/p003_s07-02.tex}
\begin{BookAnswer}
从条件出发，先求回收总量，再求剩余质量，最后求平均每批质量。\par 从问题出发，要求每批质量，需要知道剩余质量和批数；要求剩余质量，需要知道回收总量和已处理质量；要求回收总量，需要知道每班回收质量和班数。\par 两种分析方向都列同一个综合算式：\[\begin{aligned}&(4\times18-32)\div4\\={}&(72-32)\div4\\={}&40\div4\\={}&10\text{(千克)}\end{aligned}\]答：平均每批送\(10\)千克。\par 简便提示（学完运算律后）：将总量和已处理量分别平均分成4份，再求每份之差。利用乘除运算的性质，先算\(4\div4\)，同时算\(32\div4\)，避免先求72。\[\begin{aligned}&4\times18\div4-32\div4\\={}&(4\div4)\times18-32\div4\\={}&1\times18-8\\={}&18-8\\={}&10\text{(千克)}\end{aligned}\]\par 检验：\[\begin{aligned}&10\times4+32\\={}&40+32\\={}&72\text{(千克)}\end{aligned}\]与回收总量\(4\times18=72\)千克一致。
\BookFigure{figures/p003_s07-03.tex}
\BookFigure{figures/p003_s07-04.tex}
\end{BookAnswer}

\Needspace{8\baselineskip}

\pdfbookmark[2]{原书第4页}{orig-4}

\subsection*{原书第4页}

% p004_s01
\BookQuestion{四、2}
人工分拣员平均每分钟分拣40件快递，全自动分拣机每分钟分拣220件。智能机器人每分钟分拣数量比人工的3倍多20件。全自动分拣机每分钟比智能机器人多分拣多少件？\par
\begin{BookAnswer}
数量关系：机器人每分钟分拣量\(=\)人工每分钟分拣量的\(3\)倍\(+20\)；所求差\(=\)全自动分拣机每分钟分拣量\(-\)机器人每分钟分拣量。\par \par\noindent\begin{minipage}{\linewidth}\textbf{方法一：先求机器人每分钟分拣量}\par \[\begin{aligned}&220-(40\times3+20)\\&=220-(120+20)\\&=220-140\\&=80\text{(件)}\end{aligned}\]答：全自动分拣机每分钟比分拣机器人多分拣\(80\)件。\end{minipage}\par \par \par\noindent\begin{minipage}{\linewidth}\textbf{方法二：先减人工的三倍，再减多出的量}\par \[\begin{aligned}&220-40\times3-20\\&=220-120-20\\&=100-20\\&=80\text{(件)}\end{aligned}\]答：全自动分拣机每分钟比分拣机器人多分拣\(80\)件。\end{minipage}\par
\end{BookAnswer}

% p004_s02
\BookQuestion{四、3}
10月15日，学校组织270名学生和15名老师去科技馆研学。用9000元购票后还剩多少元？科技馆门票收费标准见表。旺季：4月1日～10月31日；淡季：11月1日～3月31日。\par
\par\smallskip\Needspace{4\baselineskip}\begingroup\small\setlength{\tabcolsep}{3pt}
\noindent\begin{tabularx}{\linewidth}{|*{3}{>{\centering\arraybackslash}X|}}\hline
季节 & 成人 & 儿童\\\hline
旺季 & 45元 & 30元\\\hline
淡季 & 25元 & 15元\\\hline
\end{tabularx}\par
\endgroup\smallskip
\begin{BookAnswer}
\(10\)月\(15\)日在旺季，学生票每人\(30\)元，教师票每人\(45\)元。\par \par\noindent\begin{minipage}{\linewidth}\textbf{方法一：分别求两类票价，再算剩款}\par \[\begin{aligned}&9000-(30\times270+45\times15)\\&=9000-(8100+675)\\&=9000-8775\\&=225\text{(元)}\end{aligned}\]答：购票后还剩\(225\)元。\end{minipage}\par \par \par\noindent\begin{minipage}{\linewidth}\textbf{方法二：全部按学生票计，再补教师差价}\par \[\begin{aligned}&9000-[30\times(270+15)+(45-30)\times15]\\&=9000-[30\times285+15\times15]\\&=9000-[8550+225]\\&=9000-8775\\&=225\text{(元)}\end{aligned}\]答：购票后还剩\(225\)元。\end{minipage}\par \par 按“总价\(=\)单价\(\times\)数量”分别核算两类票。检验\(30\times270+45\times15+225=9000\)，票款与剩款合计等于付款。\par\textbf{简便提示（学完运算律后）}\par \[\begin{aligned}&9000-(30\times270+45\times15)\\={}&9000-(8100+675)\\={}&9000-8100-675\\={}&900-675\\={}&225\text{(元)}\end{aligned}\]
\end{BookAnswer}

% p004_s03
\BookQuestion{四、4（方案一）}
选择合适的信息，提一个三步计算解决的问题并解答。①欣欣小学买了3箱糖果。②每箱糖果有16袋。③每箱糖果有480块。④平均分给6个年级。⑤每个年级有4个班。选择的信息：\BookBlank{}（填序号）；提出的问题：\BookBlank{}；列式解答。\par
\begin{BookAnswer}
所求是“每班的袋数”，用总袋数除以总班数；若先按年级分，再按班分，要连续平均分两次。每箱\(480\)块这一信息本问用不到。\par 选①②④⑤；问题：平均每个班分到多少袋糖果？\par \par\noindent\begin{minipage}{\linewidth}\textbf{方法一：先分年级，再分班}\par \[\begin{aligned}&16\times3\div6\div4\\&=48\div6\div4\\&=8\div4\\&=2\text{(袋)}\end{aligned}\]答：每班分到\(2\)袋。\end{minipage}\par \par \par\noindent\begin{minipage}{\linewidth}\textbf{方法二：先求总班数}\par \[\begin{aligned}&16\times3\div(4\times6)\\&=48\div24\\&=2\text{(袋)}\end{aligned}\]答：每班分到\(2\)袋。\end{minipage}\par \par\textbf{简便提示（学完运算律后）}\par \[\begin{aligned}&16\times3\div(4\times6)\\={}&16\times3\div24\\={}&16\times3\div(3\times8)\\={}&16\times3\div3\div8\\={}&16\div8\\={}&2\text{(袋)}\end{aligned}\]
\end{BookAnswer}

% p004_s04
\BookQuestion{四、4（方案二）}
选择合适的信息，提一个三步计算解决的问题并解答。①欣欣小学买了3箱糖果。②每箱糖果有16袋。③每箱糖果有480块。④平均分给6个年级。⑤每个年级有4个班。选择的信息：\BookBlank{}（填序号）；提出的问题：\BookBlank{}；列式解答。\par
\begin{BookAnswer}
本问把所求单位换成“块”：总块数\(=480\times3\)，总班数\(=4\times6\)；每班块数\(=\)总块数\(\div\)总班数。每箱\(16\)袋这一信息不需要。\par 选①③④⑤；问题：平均每个班分到多少块糖果？\par \par\noindent\begin{minipage}{\linewidth}\textbf{方法一：先分年级，再分班}\par \[\begin{aligned}&480\times3\div6\div4\\&=1440\div6\div4\\&=240\div4\\&=60\text{(块)}\end{aligned}\]答：每班分到\(60\)块。\end{minipage}\par \par \par\noindent\begin{minipage}{\linewidth}\textbf{方法二：先求总班数}\par \[\begin{aligned}&480\times3\div(4\times6)\\&=1440\div24\\&=60\text{(块)}\end{aligned}\]答：每班分到\(60\)块。\end{minipage}\par \par\textbf{简便提示（学完运算律后）}\par \[\begin{aligned}&480\times3\div(4\times6)\\={}&480\times3\div24\\={}&480\times3\div(3\times8)\\={}&480\times3\div3\div8\\={}&480\div8\\={}&60\text{(块)}\end{aligned}\]
\end{BookAnswer}

% p004_s05
\BookQuestion{五、24点（五种结构）}
“扑克牌24点”需要用四张牌上的数字通过加减乘除计算出24。下面4张牌数字是1、3、6、9。经过怎样运算可以得到24？请列出一个综合算式（\hspace{1.8em}）。\par
\BookFigure{figures/p004_s05-01.tex}
\begin{BookAnswer}
以下列出小学阶段中间结果不出现负数的五种常用结构。四个数\(1,3,6,9\)各用一次；只交换加数、因数的次序不重复列为新方法。\[\begin{aligned}&3\times(9+1)-6\\&=3\times10-6\\&=30-6\\&=24\end{aligned}\]\[\begin{aligned}&9+3\times(6-1)\\&=9+3\times5\\&=9+15\\&=24\end{aligned}\]\[\begin{aligned}&(9-1)\times(6-3)\\&=8\times3\\&=24\end{aligned}\]\[\begin{aligned}&6+9\times(3-1)\\&=6+9\times2\\&=6+18\\&=24\end{aligned}\]\[\begin{aligned}&6\times(1+9\div3)\\&=6\times(1+3)\\&=6\times4\\&=24\end{aligned}\]\par 枚举范围：四个数各用一次，只用加、减、乘、除，每一步的中间结果均为非负整数；把仅交换加数或因数次序的重复式合并后，此范围内共有上述五种结构。
\end{BookAnswer}

\Needspace{12\baselineskip}

\section{第一单元素养评价（B）}

\Needspace{8\baselineskip}

\pdfbookmark[2]{原书第5页}{orig-5}

\subsection*{原书第5页}

% p005_s01
\BookQuestion{一、1}
计算\(12\times8+15\div3\)时，可以同时先算（\hspace{1.8em}）法和（\hspace{1.8em}）法，再算（\hspace{1.8em}）法，结果是（\hspace{1.8em}）。\par
\begin{BookAnswer}
乘法\(12\times8\)和除法\(15\div3\)分别处在两个独立部分，可以先同时计算，再相加。四空依次填“乘、除、加、\(101\)”；\allowbreak{}\(12\times8+15\div3=96+5=101\)。
\end{BookAnswer}

% p005_s02
\BookQuestion{一、2（1）}
按要求在算式\(32\times400-200\div25\)中添上括号，改变运算顺序。（不改变运算符号）（1）先减再除最后乘：\BookBlank{}。\par
\begin{BookAnswer}
要先减、再除、最后乘，先用\((400-200)\)，再把商\((400-200)\div25\)整体乘\(32\)。填\(32\times[(400-200)\div25]\)。
\end{BookAnswer}

% p005_s03
\BookQuestion{一、2（2）}
按要求在算式\(32\times400-200\div25\)中添上括号，改变运算顺序。（不改变运算符号）（2）先减再乘最后除：\BookBlank{}。\par
\begin{BookAnswer}
只把减法提前，乘除仍从左到右进行。填\(32\times(400-200)\div25\)。
\end{BookAnswer}

% p005_s04
\BookQuestion{一、3}
烟台海产品特别有名。扇贝每千克12元，梭子蟹每千克96元，即食海鲜礼盒每盒的价格比扇贝和梭子蟹单价和的3倍多10元。即食海鲜礼盒每盒（\hspace{1.8em}）元。根据题意进行分析，将流程图填写完整。\par
\BookFigure{figures/p005_s04-01.tex}
\begin{BookAnswer}
从条件开始分析：先求两种单价之和，再求这个和的3倍，最后加多出的10元。\par 方法一：\[\begin{aligned}&(12+96)\times3+10\\={}&108\times3+10\\={}&324+10\\={}&334\text{(元)}\end{aligned}\]答：每盒即食海鲜礼盒334元。\par 方法二（乘法分配律）：\[\begin{aligned}&12\times3+96\times3+10\\={}&36+288+10\\={}&324+10\\={}&334\text{(元)}\end{aligned}\]答：每盒即食海鲜礼盒334元。\par 检验：\[\begin{aligned}&(334-10)\div3\\={}&324\div3\\={}&108\end{aligned}\]与两种单价之和\(12+96=108\)一致。
\BookFigure{figures/p005_s04-02.tex}
\end{BookAnswer}

% p005_s05
\BookQuestion{一、4}
630减去78的差乘14再除以7，商是（\hspace{1.8em}）。\par
\begin{BookAnswer}
“差乘\(14\)再除以\(7\)”应列\((630-78)\times14\div7\)。\[\begin{aligned}&(630-78)\times14\div7\\&=552\times14\div7\\&=7728\div7\\&=1104\end{aligned}\]也可先把\(\times14\div7\)合并成\(\times2\)，得\(552\times2=1104\)。\par\textbf{简便提示（学完运算律后）}\par 除以7只作用在前面整个乘积上，可先用14除以7得到2；括号内的减法与14除以7互不依赖，可同一步算完。\[\begin{aligned}&(630-78)\times14\div7\\={}&(630-78)\times(14\div7)\\={}&552\times2\\={}&1104\end{aligned}\]
\end{BookAnswer}

% p005_s06
\BookQuestion{一、5①}
对比算式，在圆圈里填上“\(>\)”“\(<\)”或“\(=\)”：①\(150-75\div5\times3\ \bigcirc\ (150-75)\div5\times3\)。\par
\begin{BookAnswer}
左式\(150-75\div5\times3=150-45=105\)，右式\((150-75)\div5\times3=75\div5\times3=45\)。\(105>45\)，填\(>\)。
\end{BookAnswer}

% p005_s07
\BookQuestion{一、5②及发现}
对比算式，在圆圈里填上“\(>\)”“\(<\)”或“\(=\)”：②\(240\div(6+2)\times4\ \bigcirc\ 240\div[(6+2)\times4]\)。对比两组算式，我发现：括号能\BookBlank{}，从而改变算式的计算结果。\par
\begin{BookAnswer}
先算\(6+2=8\)。左式为\(240\div8\times4=30\times4=120\)，比30大；右式为\(240\div32\)，因为\(32>8\)，其商比\(240\div8=30\)小。因此左式大于右式，填\(>\)。\par 两组算式说明：括号能改变运算顺序，从而改变计算结果。比较时不必去求尚未学习的一般小数除法。
\end{BookAnswer}

% p005_s08
\BookQuestion{一、6①}
根据分步算式，写出综合算式：①\(36\div4=9\)，\allowbreak{}\(12+9=21\)，\allowbreak{}\(21\times5=105\)。综合算式：\BookBlank{}。\par
\begin{BookAnswer}
把\(36\div4\)代入\(12+9\)，再把和代入\(21\times5\)。填\((12+36\div4)\times5\)。
\end{BookAnswer}

% p005_s09
\BookQuestion{一、6②}
根据分步算式，写出综合算式：②\(60-35=25\)，\allowbreak{}\(25\times4=100\)，\allowbreak{}\(500\div100=5\)。综合算式：\BookBlank{}。\par
\begin{BookAnswer}
先算\(60-35\)，再乘\(4\)，这个积作除数。填\(500\div[(60-35)\times4]\)。
\end{BookAnswer}

\Needspace{8\baselineskip}

\pdfbookmark[2]{原书第6页}{orig-6}

\subsection*{原书第6页}

% p006_s01
\BookQuestion{一、7}
导航显示，从家到公园全程1800米，小明先以每分钟80米的速度走了5分钟，剩下的路程要在15分钟内走完，剩下的路程每分钟需要走多少米？综合算式：\BookBlank{}。\par
\begin{BookAnswer}
先求已走路程\(80\times5=400\)米，再求剩余路程\(1800-400=1400\)米。\par 所求平均每分钟走的米数\(=\)剩余路程\(\div\)剩余时间，所以原空填\((1800-80\times5)\div15\)。本题要求列综合算式，到此即可。
\end{BookAnswer}

% p006_s02
\BookQuestion{一、8}
学校体育室准备给同学们发放运动器材。①一共有15盒羽毛球。②每盒12个。③已经分给同学们80个。④剩下的羽毛球要分给5个班级。平均每个班分几个？通过信息①和信息②可以求出（\hspace{1.8em}），然后与信息③合并可以求出（\hspace{1.8em}），最后再求出剩下的羽毛球平均每个班级能分到（\hspace{1.8em}）个。\par
\begin{BookAnswer}
文字空依次填“一共有多少个羽毛球”“剩下多少个羽毛球”；数值空填\(20\)。\par \par\noindent\begin{minipage}{\linewidth}\textbf{方法一：先求总数，再扣除已分数量}\par \[\begin{aligned}&(12\times15-80)\div5\\&=(180-80)\div5\\&=100\div5\\&=20\text{(个)}\end{aligned}\]答：平均每班分到\(20\)个羽毛球。\end{minipage}\par \par \par\noindent\begin{minipage}{\linewidth}\textbf{方法二：总数和已分数量分别平均分}\par \[\begin{aligned}&12\times15\div5-80\div5\\&=180\div5-16\\&=36-16\\&=20\text{(个)}\end{aligned}\]答：平均每班分到\(20\)个羽毛球。\end{minipage}\par \par 检验\(20\times5+80=180\)，且\(12\times15=180\)：分给各班的数量加上已经分发的数量，正好等于原有总数。\par\textbf{简便提示（学完运算律后）}\par \[\begin{aligned}&12\times15\div5-80\div5\\={}&12\times(15\div5)-80\div5\\={}&12\times3-16\\={}&36-16\\={}&20\text{(个)}\end{aligned}\]
\end{BookAnswer}

% p006_s03
\BookQuestion{一、9}
\(\square+\square+\square+\square=100\)，\allowbreak{}\((\triangle+\triangle)\times\square=100\)，\allowbreak{}\(\square\times\triangle\times\triangle=\)（\hspace{1.8em}）。\par
\begin{BookAnswer}
四个\(\square\)相加是\(100\)，故\(\square=100\div4=25\)。再由\((\triangle+\triangle)\times25=100\)，得\(\triangle+\triangle=4\)，所以\(\triangle=2\)。所求\(25\times2\times2=100\)，填\(100\)。\par\textbf{简便提示（学完运算律后）}\par \[\begin{aligned}&25\times2\times2\\={}&25\times(2\times2)\\={}&25\times4\\={}&100\end{aligned}\]
\end{BookAnswer}

% p006_s04
\BookQuestion{二、1}
给算式\(600\div24+6\times2\)添加括号，使运算顺序变为“先求和，再求积，最后求商”，下列添加方式正确的是（\hspace{1.8em}）。\par
\noindent A.\(600\div24+(6\times2)\)\par
\noindent B.\(600\div(24+6)\times2\)\par
\noindent C.\(600\div[(24+6)\times2]\)\par
\noindent D.\((600\div24+6)\times2\)\par
\begin{BookAnswer}
选\(\mathrm C\)。先和\(24+6\)，再积\((24+6)\times2\)，最后用\(600\)除以这个积，得\(600\div[(24+6)\times2]\)。
\end{BookAnswer}

% p006_s05
\BookQuestion{二、2}
将“\(520+360=880\)”，“\(900-880=20\)”，“\(160\div20=8\)”三个算式合并成一个综合算式，正确的是（\hspace{1.8em}）。\par
\noindent A.\(900-(520+360)\div20\)\par
\noindent B.\(160\div[900-(520+360)]\)\par
\noindent C.\(160\div(900-520+360)\)\par
\noindent D.\(160\div(900-520-360)\)\par
\begin{BookAnswer}
按给出的三个分步算式逐层合成，选\(\mathrm B\)：\(160\div[900-(520+360)]\)。\par 注意：\(\mathrm D\)的\(160\div(900-520-360)\)与\(\mathrm B\)数值相等，但它没有按原先“先把\(520\)与\(360\)相加”的结构合成。若题目只要求结果相等，\allowbreak{}\(\mathrm B,\mathrm D\)都得\(8\)；本题按给定步骤合并，采用\(\mathrm B\)。
\end{BookAnswer}

% p006_s06
\BookQuestion{二、3（两个空）}
\(25\times8\div25\times8=\)（\hspace{1.8em}），\allowbreak{}\(25\times8\div(25\times8)=\)（\hspace{1.8em}）。\par
\noindent A.1\par
\noindent B.64\par
\noindent C.0\par
\noindent D.625\par
\begin{BookAnswer}
第一空选\(\mathrm B\)：\(25\times8\div25\times8=200\div25\times8=64\)。\par 第二空选\(\mathrm A\)：\(25\times8\div(25\times8)=200\div200=1\)。括号使两个乘积整体相除。\par\textbf{简便提示（学完运算律后）}\par \[\begin{aligned}&25\times8\div25\times8\\={}&(25\div25)\times(8\times8)\\={}&1\times64\\={}&64\end{aligned}\]
\end{BookAnswer}

% p006_s07
\BookQuestion{二、4}
被减数、减数与差的和是120，那么被减数是（\hspace{1.8em}）。\par
\noindent A.120\par
\noindent B.60\par
\noindent C.40\par
\noindent D.条件不够\par
\begin{BookAnswer}
选\(\mathrm B\)，被减数是\(60\)。因为“减数\(+\)差\(=\)被减数”，三个数的和等于被减数的\(2\)倍，所以\(120\div2=60\)。
\end{BookAnswer}

% p006_s08
\BookQuestion{二、5}
要使\(8\times(\square\div3)\div2=24\)，下列说法和计算结果正确的是（\hspace{1.8em}）。\par
\noindent A.先算\(8\div2=4\)，再算\(4\times24=96\)，最后算\(96\div3=32\)\par
\noindent B.先算\(24\times2=48\)，再算\(48\div8=6\)，最后算\(6\times3=18\)\par
\noindent C.先算\(24\times2=48\)，再算\(48\div8=6\)，最后算\(6\div3=2\)\par
\noindent D.先算\(8\div2=4\)，再算\(24\div4=6\)，最后算\(6\div3=2\)\par
\begin{BookAnswer}
选\(\mathrm B\)。从结果倒推：先\(24\times2=48\)，再\(48\div8=6\)，最后\(6\times3=18\)。检验：\(8\times(18\div3)\div2=24\)。
\end{BookAnswer}

\Needspace{8\baselineskip}

\pdfbookmark[2]{原书第7页}{orig-7}

\subsection*{原书第7页}

% p007_s01
\BookQuestion{二、6}
某农田去年产小麦6800千克，今年比去年多1200千克，明年计划产量比今年的2倍少500千克，明年计划产小麦多少千克？正确算式是（\hspace{1.8em}）。\par
\noindent A.\((6800+1200)\times2-500\)\par
\noindent B.\(6800+1200\times2-500\)\par
\noindent C.\(6800\times2+1200-500\)\par
\noindent D.\((6800+1200\times2)-500\)\par
\begin{BookAnswer}
选\(\mathrm A\)。先求今年产量\(6800+1200\)，再求今年的\(2\)倍，最后减\(500\)。列式\((6800+1200)\times2-500\)，结果是\(15500\)千克。
\end{BookAnswer}

% p007_s02
\BookQuestion{三、1}
计算下列各题：\(3744-[(85-47)\times6]\)。\par
\begin{BookAnswer}
\[\begin{aligned}&3744-[(85-47)\times6]\\&=3744-[38\times6]\\&=3744-228\\&=3516\end{aligned}\]
\end{BookAnswer}

% p007_s03
\BookQuestion{三、2}
计算下列各题：\(25\times[(160+512)\div32]\)。\par
\begin{BookAnswer}
\[\begin{aligned}&25\times[(160+512)\div32]\\&=25\times[672\div32]\\&=25\times21\\&=525\end{aligned}\]\par\textbf{简便提示（学完运算律后）}\par \[\begin{aligned}&25\times21\\={}&25\times(20+1)\\={}&500+25\\={}&525\end{aligned}\]
\end{BookAnswer}

% p007_s04
\BookQuestion{三、3}
计算下列各题：\([275\div(23+32)]\times45\)。\par
\begin{BookAnswer}
\[\begin{aligned}&[275\div(23+32)]\times45\\&=[275\div55]\times45\\&=5\times45\\&=225\end{aligned}\]
\end{BookAnswer}

% p007_s05
\BookQuestion{三、4}
计算下列各题：\([73\times(38+55)]-2876\)。\par
\begin{BookAnswer}
\[\begin{aligned}&[73\times(38+55)]-2876\\&=[73\times93]-2876\\&=6789-2876\\&=3913\end{aligned}\]\par\textbf{简便提示（学完运算律后）}\par 括号按原顺序算得93后，把93看成100少7。先多算7份73再减回；乘100与乘7在同一步完成。\[\begin{aligned}&73\times93\\={}&73\times(100-7)\\={}&73\times100-73\times7\\={}&7300-511\\={}&6789\end{aligned}\]
\end{BookAnswer}

% p007_s06
\BookQuestion{三、5}
计算下列各题：\(65\times[(384+576)\div12]\)。\par
\begin{BookAnswer}
\[\begin{aligned}&65\times[(384+576)\div12]\\&=65\times[960\div12]\\&=65\times80\\&=5200\end{aligned}\]
\end{BookAnswer}

% p007_s07
\BookQuestion{三、6}
计算下列各题：\(336\div[252\div(12\times3)]\)。\par
\begin{BookAnswer}
\[\begin{aligned}&336\div[252\div(12\times3)]\\&=336\div[252\div36]\\&=336\div7\\&=48\end{aligned}\]
\end{BookAnswer}

% p007_s08
\BookQuestion{四、1}
学校“爱心实践社”为乡村小学“梦想书屋”采购阅览桌和配套椅子，他们选了2张阅览桌和5把椅子，电子小票如表。请算一算每把椅子的单价是多少元。（先分析，再解答。）\par
\par\smallskip\Needspace{4\baselineskip}\begingroup\small\setlength{\tabcolsep}{3pt}
\noindent\textbf{购物小票}\par\smallskip
\noindent\begin{tabularx}{\linewidth}{|*{5}{>{\centering\arraybackslash}X|}}\hline
编号 & 商品名称 & 单位 & 数量 & 单价\\\hline
1 & 桌子 & 张 & 2 & 80元\\\hline
2 & 椅子 & 把 & 5 & 遮损未显示\\\hline
\end{tabularx}\par
\noindent 总计金额：230元\par
\endgroup\smallskip
\BookFigure{figures/p007_s08-01.tex}
\begin{BookAnswer}
从问题开始分析：椅子单价需要椅子总价和椅子数量；椅子总价需要总金额和桌子总价；桌子总价需要桌子单价和桌子数量。\par 方法一：\[\begin{aligned}&(230-80\times2)\div5\\={}&(230-160)\div5\\={}&70\div5\\={}&14\text{(元)}\end{aligned}\]答：每把椅子的单价是14元。\par 方法二：将总金额与桌款分别平均分成5份，再求每份的差。\[\begin{aligned}&230\div5-80\times2\div5\\={}&46-160\div5\\={}&46-32\\={}&14\text{(元)}\end{aligned}\]答：每把椅子的单价是14元。\par 检验：\[\begin{aligned}&80\times2+14\times5\\={}&160+70\\={}&230\text{(元)}\end{aligned}\]与小票总金额一致。
\BookFigure{figures/p007_s08-02.tex}
\end{BookAnswer}

\Needspace{8\baselineskip}

\pdfbookmark[2]{原书第8页}{orig-8}

\subsection*{原书第8页}

% p008_s01
\BookQuestion{四、2}
近年来，“共享菜园”逐渐流行，我市推出了“共享菜园”租赁方案：每块菜地每月租金66元，满2000元减300元。张叔叔租了5块这样的菜地并计划使用8个月，张叔叔实际需要付多少钱？\par
\BookFigure{figures/p008_s01-01.tex}
\begin{BookAnswer}
优惠前租金\(=\)每块地每月租金\(\times\)块数\(\times\)月数。下列计算得到原租金\(2640\)元，\allowbreak{}\(2640>2000\)，满足优惠条件，才减\(300\)元。\par \par\noindent\begin{minipage}{\linewidth}\textbf{方法一：先求五块地每月的租金}\par \[\begin{aligned}&66\times5\times8-300\\&=330\times8-300\\&=2640-300\\&=2340\text{(元)}\end{aligned}\]答：张叔叔实际需要付\(2340\)元。\end{minipage}\par \par \par\noindent\begin{minipage}{\linewidth}\textbf{方法二：先求每块地八个月的租金}\par \[\begin{aligned}&66\times8\times5-300\\&=528\times5-300\\&=2640-300\\&=2340\text{(元)}\end{aligned}\]答：张叔叔实际需要付\(2340\)元。\end{minipage}\par \par\textbf{简便提示（学完运算律后）}\par \[\begin{aligned}&66\times5\times8-300\\={}&66\times(5\times8)-300\\={}&66\times40-300\\={}&2640-300\\={}&2340\text{(元)}\end{aligned}\]
\end{BookAnswer}

% p008_s02
\BookQuestion{四、3}
烟台海阳海上火箭发射闻名全国。某天海阳航天科普馆上午接待了420名学生，下午参观的学生人数比上午的2倍少60人。下午参观的学生比上午多多少人？\par
\begin{BookAnswer}
下午人数\(=\)上午人数\(\times2-60\)；题目求的是两者之差，还要减去上午人数。求出差后可用\(420+360=780\)核对下午人数。\par \par\noindent\begin{minipage}{\linewidth}\textbf{方法一：先算下午人数，再比较}\par \[\begin{aligned}&420\times2-60-420\\&=840-60-420\\&=780-420\\&=360\text{(人)}\end{aligned}\]答：下午比上午多\(360\)人。\end{minipage}\par \par \par\noindent\begin{minipage}{\linewidth}\textbf{方法二：直接利用两倍关系求差}\par \[\begin{aligned}&420\times(2-1)-60\\&=420\times1-60\\&=420-60\\&=360\text{(人)}\end{aligned}\]答：下午比上午多\(360\)人。\end{minipage}\par 下午的“两份上午人数”减去一份上午人数，还剩一份，再减去\(60\)人即可。
\end{BookAnswer}

% p008_s03
\BookQuestion{四、4}
考古发现一段战国城墙，长440米。其中一部分已经修复，修复的长度是未修复的3倍，未修复的部分计划用5天修完，平均每天修多少米？\par
\begin{BookAnswer}
未修复的长度是\(1\)份，修复的是\(3\)份，总长是\(4\)份。\par \par\noindent\begin{minipage}{\linewidth}\textbf{方法一：先求未修复长度}\par \[\begin{aligned}&440\div(3+1)\div5\\&=440\div4\div5\\&=110\div5\\&=22\text{(米)}\end{aligned}\]答：平均每天须修\(22\)米。\end{minipage}\par \par \par\noindent\begin{minipage}{\linewidth}\textbf{方法二：合并两次平均分}\par \[\begin{aligned}&440\div[(3+1)\times5]\\&=440\div[4\times5]\\&=440\div20\\&=22\text{(米)}\end{aligned}\]答：平均每天须修\(22\)米。\end{minipage}\par \par \par\noindent\begin{minipage}{\linewidth}\textbf{方法三：先求已修复长度再扣除}\par \[\begin{aligned}&[440-440\div(3+1)\times3]\div5\\&=[440-440\div4\times3]\div5\\&=[440-110\times3]\div5\\&=[440-330]\div5\\&=110\div5\\&=22\text{(米)}\end{aligned}\]答：平均每天须修\(22\)米。\end{minipage}\par \par 检验：\(22\times5=110\)米是未修部分，\allowbreak{}\(110\times3=330\)米是已修部分，\allowbreak{}\(110+330=440\)米，份数关系和总长都符合条件。\par\textbf{简便提示（学完运算律后）}\par \[\begin{aligned}&440\div4\div5\\={}&440\div[4\times5]\\={}&440\div20\\={}&22\text{(米)}\end{aligned}\]
\end{BookAnswer}

% p008_s04
\BookQuestion{五、1}
甲数除以乙数的商是5，甲、乙两数的和是300，甲数是（\hspace{1.8em}），乙数是（\hspace{1.8em}）。\par
\begin{BookAnswer}
乙数是\(1\)份，甲数是\(5\)份，总和\(300\)是\(6\)份。\(300\div(5+1)=50\)，所以乙数\(50\)，甲数\(50\times5=250\)。按题目两个空的顺序填\(250,50\)。
\end{BookAnswer}

% p008_s05
\BookQuestion{五、2（括号条件）}
给\(1200\div30-20\div2\)添上小括号，结果最大是（\hspace{1.8em}），算式\BookBlank{}；结果最小是（\hspace{1.8em}），算式\BookBlank{}。\par
\begin{BookAnswer}
题面没有写“只添一对”。\textbf{允许嵌套小括号时，最大值是\(240\)}：\[\begin{aligned}&1200\div((30-20)\div2)\\&=1200\div(10\div2)\\&=1200\div5\\&=240\end{aligned}\]\textbf{如果课堂另限只添一对小括号，最大值才是\(60\)}，两种写法为\(1200\div(30-20)\div2\)与\(1200\div(30-20\div2)\)。两种范围下的\textbf{最小值都是\(10\)}：\[\begin{aligned}&(1200\div30-20)\div2\\&=(40-20)\div2\\&=20\div2\\&=10\end{aligned}\]保持数与运算符顺序不变，五种完全加括号结构为：\par \(((1200\div30)-20)\div2=10\)；\par \((1200\div(30-20))\div2=60\)；\par \((1200\div30)-(20\div2)=30\)；\par \(1200\div((30-20)\div2)=240\)；\par \(1200\div(30-(20\div2))=60\)。\par 因此以上极值已包含所有结构。
\end{BookAnswer}
